\documentclass[12pt]{article}
\usepackage[utf8]{inputenc}
\usepackage[T1]{fontenc}
\usepackage{hyperref}
\title{Taward\^e Language Resources for K-12 Curriculum Development}
\author{Compiled by Claude Sonnet 5.5 (claude-sonnet-5-5) for Andr\'as Kornai}
\date{}
\begin{document}
\maketitle

\section{Introduction}
Taward\^e (ISO 639-3 xtw; Glottocode tawa1278) is a Northern Nambikwaran language of Mato Grosso, Brazil. Other names include Nambiquara do Norte and Da'wan'du. Sources disagree on whether it is a separate language or a dialect within a single Nambikwara language. I found no Taward\^e-specific speaker count. The whole family is estimated at about 1,000 people. Joshua Project's page for the Tawande reports a population of 20, but it lists the primary language as Sabane (sae), so it should be treated with caution for this language.

\section{Primary Educational Resources}
No primers, textbooks or school programmes specific to Taward\^e were found.

\section{Linguistic and Reference Materials}
\begin{itemize}
\item Wikipedia, Nambikwaran languages: \url{https://en.wikipedia.org/wiki/Nambikwaran_languages}
\item Wikipedia, Nambikwara language: \url{https://en.wikipedia.org/wiki/Nambikwara_language}
\item PanLex entry for ISO 639-3 xtw: \url{https://apps.panlex.org/panlinx/ex/2746683}
\end{itemize}

\section{Religious and Community Resources}
\begin{itemize}
\item Joshua Project, Tawande in Brazil (see the caution above about the language listed): \url{https://joshuaproject.net/people_groups/19354}
\end{itemize}

\section{Digital and Audio Resources}
No audio, video or corpus specific to Taward\^e was found.

\section{Curriculum Development Context}
Little is documented and the language is said to have only elderly speakers. Any work would begin with community-led documentation. Materials from other Nambikwaran languages, such as Mamaind\^e, might be a starting point, but this search did not examine them.
\end{document}